\note{1}{Fri 23 Feb 2024 14:42}{Complex Analysis Review}
\tags{Winding Number}

\subsection{Cauchy's Theorem}

\subsection{The Winding Number}

We consider the following integral:
\[
	\oint_{S^1} \frac{dz }{z} = 2 \pi i
.\] 

\begin{intuition}
	As the path traverses through $S^1$, $dz $ and $z$ are orthogonal, so $dz = i z dx$ for some (hyper-)real number $dx$. Hence, $\frac{dz}{z} = i dx$. What is this real number $dx$? From geometry we can interpret it as the angle swiped by arc length $dz$. Thus,
	\[
		\oint_{S^1} \frac{dz }{z} = \int _{0}^{2 \pi} i d \theta = 2 \pi i
	.\] 
\end{intuition}
\begin{proof}
	Parametrize with the exponential, which is an analytic function. We have
	\[
		\int _{S^1} \frac{dz }{z}
		= \int _0^{2 \pi} \frac{d(e^{i t})}{e^{i t}}
		= \int _0^{2 \pi} \frac{i e^{i t}}{e^{i t}} = 2 \pi i
	.\] 
\end{proof}

\subsection{Cauchy Integral Formula}

Two key ingredients:
\begin{itemize}
	\item As a direct consequence, all holomorphic functions (locally differentiable) are analytic (infinitely differentiable and has series expansion).
	\item If $D \subset \mathbb{C}$ is a simple closed region with trivial topology then the value of $f$ in $\text{int} D$ is determined by its value on $\partial D$. In particular, if $f$ is not analytic at some point in $D$ but has an analytic continuation, the formula prescribes that continuation.
\end{itemize}






